\section{Solución Propuesta}
Se propone el desarrollo de una plataforma web y móvil integral que actúe como un ecosistema transaccional y de geolocalización. Para asegurar la factibilidad técnica y respetar las restricciones de recursos humanos y temporales de la carrera, el alcance del Producto Mínimo Viable (MVP) se estructurará de forma modular, priorizando la carga manual y la reserva programada.

\subsection{División en Módulos de Trabajo}
\begin{itemize}
    \item \textbf{Módulo de Mapeo y Geolocalización:} Visualización de plazas disponibles sobre mapas digitales interactivos mediante la integración de APIs de cartografía.
    \item \textbf{Módulo de Reserva Dinámica:} Motor lógico para la gestión de disponibilidad horaria y reservas manuales en tiempo real.
    \item \textbf{Módulo de Autenticación y Seguridad:} Gestión de perfiles diferenciados para conductores (demandantes) y propietarios/playas (oferentes).
    \item \textbf{Módulo de Pagos:} Integración con procesadores de pago y billeteras virtuales para la retención automática de comisiones y liquidación a proveedores.
\end{itemize}

\subsection{Nota sobre el Alcance y Restricciones}
Si bien la visión a largo plazo contempla la integración de sensores IoT para la detección automática de ocupación física en tiempo real y componentes de Inteligencia Artificial descriptiva para predicción de demanda, se concluye que dichos elementos representan un riesgo de complejidad elevado para la fase inicial. Por lo tanto, se posponen como evoluciones futuras del sistema.